\documentclass[10pt,a4paper]{amsart}
\usepackage[margin=19mm]{geometry}
\usepackage{amsmath,amssymb,mathtools}
\usepackage[scr=rsfs,cal=euler]{mathalfa}
\usepackage{xfrac}
\usepackage[unicode,hidelinks,bookmarks=false]{hyperref}
\newcommand{\ie}{i.e.\ }
\newcommand{\cf}{cf.\ }
\newcommand{\resp}{respectively\ }
\newcommand{\Subsection}{\subsection}
\providecommand{\nobreakdash}{\nobreak\hbox{-}}
\setlength{\emergencystretch}{2em}
\multlinegap0pt
\usepackage{bm}
\usepackage{bbm}
\usepackage{dsfont}
\usepackage{xspace}
\usepackage{cancel}
\usepackage{mathtools}
\usepackage{amsthm}
\makeatletter
\@ifundefined{theorem}{\newtheorem{theorem}{Theorem}}{}
\makeatother
\title[Artificial Bee Colony-based Adaptive Position Control of Ele]{Artificial Bee Colony-based Adaptive Position Control of Electrohydraulic Servo Systems with Parameter Uncertainty}
\author{Babajide O. Ayinde, Sami El-Ferik}
\date{Source: arXiv:1710.03881v1 (CC BY 4.0)}
\begin{document}
\maketitle

\noindent\emph{Source and licence.} Artificial Bee Colony-based Adaptive Position Control of Electrohydraulic Servo Systems with Parameter Uncertainty. Version arXiv:1710.03881v1 (CC BY 4.0), distributed under the Creative Commons Attribution 4.0 International licence, as recorded in the arXiv licence metadata for this exact version and shown on the abstract page https://arxiv.org/abs/1710.03881v1. Full rights notice: licenses/arxiv-1710.03881-CC-BY-4.0.txt.\par\smallskip

\noindent\textit{Editorial scope.} Counted unit: the Theorem whose source label is \texttt{MyEq14} in the article (source0.tex lines 184--234, source label \texttt{MyEq14}), delivered together with the article's own complete original proof of it (source0.tex lines 235--291, one \texttt{proof} environment of 57 lines). No mathematics is written, repaired or re-derived: statement and proof are the article's own text line for line. Required context retained uncounted: MyEq4 (lines 107--113); MyEq87a (lines 161--169). Every definition, display and label of those lines is retained unchanged. Omitted from this package and not redistributed: 1--106, 114--160, 170--183, 292--431. Nothing inside a retained excerpt is omitted or altered: every retained body is delivered exactly as the article's lines are written, so no omission receipt of any kind accompanies this package. Citations: the retained excerpts carry no citation command at all, so no bibliography entry of the article is transcribed and no reference is redistributed. Reference adaptation: none was needed and none was made; every reference inside the delivered unit resolves inside it, so no unresolved reference or citation is produced. Printed slips kept literal: no printed word, symbol, hypothesis or number of the counted statement or of its proof was altered. Layout and class: the article's own class is \texttt{article} with packages including nips14submit\_e, times, moreverb, url, graphicx, amsmath, amssymb, caption; the wrapper is the lane's pinned \texttt{amsart} editorial wrapper at 10pt on A4, which restates the theorem-like environments the retained text needs and adds the article's own macro definitions that the retained text needs (none), with extra wrapper packages (bm, bbm, dsfont, xspace, cancel, mathtools). The delivered numbering is the unit's own. Assets: no retained span contains a figure environment, a graphics-inclusion command, a PDF-page inclusion, a table or a TikZ picture; no image file is copied into this package, the package declares no asset dependency and no image file is redistributed with it. Licence: version arXiv:1710.03881v1 of this article is distributed under the Creative Commons Attribution 4.0 International licence, as recorded in the arXiv licence metadata for this exact version and shown on the abstract page https://arxiv.org/abs/1710.03881v1. A full rights notice is delivered at licenses/arxiv-1710.03881-CC-BY-4.0.txt. Characters: every retained excerpt is pure ASCII, so the delivered engine, XeLaTeX with its default Latin Modern text font, reports no missing character.
 \begin{equation} \label{MyEq4}
  \begin{split}
  {\dot \xi}_1 &= \xi_2+d(t),\\
  {\dot \xi}_2 &= \tfrac{1}{m_t}(S \xi_3-b\xi_2-\beta \xi_1),\\
  {\dot \xi}_3 &= \tfrac{4B}{V_t}(ku\sqrt{P_d-\text{sign}(u)\xi_3}-\tfrac{\alpha \xi_3}{1+\gamma \lvert u\rvert}-S\xi_2 ).
  \end{split}
\end{equation}

\begin{equation} \label{MyEq87a}
  \begin{split}
  {\dot e_1} &= e_2 + d,\\
  {\dot e_2} &= e_3,\\
  {\dot e_3} &= \tfrac{b_0}{m_t^2}\theta^T\phi(\xi)\xi_1 +(-\tfrac{1}{m_t}\theta^T\phi(\xi)+\tfrac{b_0^2}{m_t^2}-\tfrac{4BS^2}{m_tV_t})\xi_2\\
  &+(-\tfrac{b_0S}{m_t^2}-\tfrac{4BS\alpha}{m_tV_t(1+\gamma \lvert u \rvert)})\xi_3 + \Delta F_1\theta^T\phi(\xi)\xi_1\\
  & +\Delta F_2\xi_2-\Delta F_3\xi_3-\theta_d^T\phi(\xi)d(t)-{\dddot r} + Am(t)u.
  \end{split}
\end{equation}

\begin{theorem}
Given that
\begin{align} \label{MyEq14}
& \alpha_4= \tfrac{3}{2}+\lambda ,\nonumber \\
& \alpha_5= 1 + \tfrac{1}{\lambda^3}+\tfrac{1}{2 \lambda},\\
& \alpha_6=\tfrac{1}{\lambda^4}\nonumber.
\end{align}
and
\begin{equation} \label{MyEq99}
  \begin{split}
   h(e) &= \alpha_4e_1+\alpha_5e_2+\alpha_6e_3\\
   g(e,\xi) &= e_2+\lambda e_1+\alpha_4e_2+\alpha_5e_3\\
   &+\alpha_6(\tfrac{b_0^2}{m_t^2}-\tfrac{4BS^2}{m_tV_t})\xi_2-\tfrac{b_0S}{m_t^2}\xi_3\\
   A_1 & =\tfrac{\alpha_6b_0}{m_t^2}, \ \ \ A_2=\tfrac{\alpha_6}{m_t} \ \ \  A_3=\tfrac{4BS\alpha\alpha_6}{m_tV_t}\\
   \tilde{\theta} & = \theta -\hat \theta
  \end{split}
\end{equation}
and let the adaptation law be given as
\begin{equation}
  \begin{split}\label{MyEq99aaa}
  {\dot {\hat \theta}} &=\gamma_6(A_1h(e)\phi(\xi)\xi_1-A_2h(e)\phi(\xi)\xi_2\\
                       &+\alpha_6F_{max}|\xi_1|\phi(\xi))\\
  {\dot {\hat \theta}_d} &=-\gamma_7\alpha_6\phi(\xi)d_{max}
  \end{split}
\end{equation}
and let the adaptive feedback be given as
\begin{equation}\label{MyEq2ccc}
  \begin{split}
  u=\bigg(\tfrac{m_tV_t}{4\alpha_3SBk \text{min}(\sqrt{P_d-\xi_3},\sqrt{P_d+\xi_3}}\bigg)v
  \end{split}
\end{equation}
where,
\begin{equation}
  \begin{split}
  v=&-(|g(e,\xi)|+|\alpha_4-\alpha_6{\hat \theta}^T_d\phi(\xi)|d_{max}+\Phi(\xi,\hat \theta)\\
    &+\alpha_6|\xi_3|F_{max}+\alpha_6|\xi_2|F^3_{max})-k_oh(e)
  \end{split}
\end{equation}
Then, system \eqref{MyEq4} under the adaptive feedback control law given in \eqref{MyEq2ccc} is practically stable and the solution of the error dynamic \eqref{MyEq87a} is globally uniformly ultimately bounded with ultimate bound satisfying the following condition
\begin{equation}
  \begin{split}
||e||^2 \leq \frac{d_{max}^2}{\lambda\,\sigma_{min}(\phi \phi^T)}\leq\frac{d_{max}^2}{2 \left(\lambda+\lambda^{5}\right)\,\sigma_{min}(\phi \phi^T)}
\end{split}
\end{equation}
with
\begin{equation} \label{phi2}
  \begin{split}
\phi = \left[ \begin{array}{lll} 1&\lambda&\frac{\alpha_1}{\lambda^2}\\ 0 & 1 & \frac{\alpha_2}{\lambda^2}\\ 0& 0&\frac{\alpha_3}{\lambda^2} \end{array}\right]
\end{split}
\end{equation}
\end{theorem}

\begin{proof}
To prove the boundedness of the error dynamics, we again choose the Lyapunov functions as follows:
\begin{equation} \label{MyEq88}
  \begin{split}
  V_1 &= \tfrac{1}{2}e_1^2,\ \ \ V_2=\tfrac{1}{2\lambda^4}(e_2+\lambda e_1)^2\\
   V_3 &=\tfrac{1}{2}(\alpha_4 e_1 + \alpha_5 e_2 + \alpha_6 e_3)^2\\ V_4 &=\tfrac{1}{2\gamma_6}{\tilde{\theta}}^T{\tilde{\theta}}^T+\tfrac{1}{2\gamma_7}{\tilde{\theta_d}}^2,
  \end{split}
\end{equation}
and again using the Young's inequality with  $\lambda >0$ then,
\begin{equation} \label{MyEq103}
  \begin{split}
  {\dot V}= {\dot V_1}+{\dot V_2}+{\dot V_3}+{\dot V_4}
  \end{split}
\end{equation}
Therefore, the derivative of the Lyapunov function is thus given:
\begin{equation} \label{MyEq104}
  \begin{split}
  {\dot V} & \leq  -\tfrac{\lambda}{2} e_1^2+ (\tfrac{1}{2\lambda}+\tfrac{1}{2\lambda^5})d^2-(e_2+ \lambda e_1)^2\\
  &+ h(e)[g(e,\xi)+A_1\tilde{\theta}^T\phi(\xi)\xi_1+A_1{\hat\theta}^T\phi(\xi)\xi_1\\
  &-A_2\tilde{\theta}^T\phi(\xi)\xi_2-A_2{\hat\theta}^T\phi(\xi)\xi_2-\tfrac{A_3}{1+\gamma \lvert u \rvert})\xi_3 \\
  &+ \alpha_4d+  \alpha_6\Delta F_1\tilde{\theta}^T\phi(\xi)\xi_1+\alpha_6\Delta F_1{\hat\theta}^T\phi(\xi)\xi_1\\
  &+\alpha_6\Delta F_2\xi_2-\alpha_6\Delta F_3\xi_3-\alpha_6\tilde{\theta}_d^T\phi(\xi)d(t)\\
  &-\alpha_6{\hat\theta}_d^T\phi(\xi)d(t)-\alpha_6{\dddot r} + \alpha_6Am(t)u]\\
   &-\tfrac{1}{\gamma_6}{\tilde{\theta}}^T{\dot {\hat \theta}}-\tfrac{1}{\gamma_7}\tilde{\theta}_d{\dot {\hat \theta}_d}
  \end{split}
\end{equation}
In order to annihilate the parametric error, the update laws are chosen as in \eqref{MyEq99aaa}
Given that $F_{max}=\sup_{t \geq 0}|\Delta F_1|$ and $d_{max}= \sup_{t \geq 0}|d(t)|$, and to ensure a uniformly ultimately bounded tracking error, adaptive feedback $``u"$ is chosen as given in \eqref{MyEq2ccc} and $v$ is given as follows:
\begin{equation} \label{MyEq110}
  \begin{split}
  v&=-(|g(e,\xi)|+|\alpha_4-\alpha_6{\hat \theta}^T_d\phi(\xi)|d_{max}+\Phi(\xi,\hat \theta)\\
  &+\alpha_6|\xi_3|F_{max}+\alpha_6|\xi_2|F^3_{max})-k_oh(e)
  \end{split}
\end{equation}
Ultimately,
\begin{equation} \label{MyEq111}
  \begin{split}
{\dot V}  &\leq -\tfrac{\lambda}{2} e_1^2+ (\tfrac{1}{2\lambda}+\tfrac{1}{2\lambda^5})d^2-(e_2+ \lambda e_1)^2\\
&-(\alpha_1 e_1 + \alpha_2 e_2 + \alpha_3 e_3)^2
  \end{split}
\end{equation}
Therefore,for ${\dot V} \leq 0$, it is sufficient to verify that $-V  +\tfrac{1}{\lambda}d_{max}^2 \leq 0$. Which means that ${\dot V} \geq 0$ if $e$ is such that $V \leq \tfrac{1}{\lambda}d_{max}^2$. This will lead to increasing $e$ until $V \geq \tfrac{1}{\lambda}d_{max}^2$.\\
Let $z=\phi^T e$, with
\begin{equation} \label{weights1}
  \begin{split}
\phi = \left[ \begin{array}{lll} \frac{1}{\sqrt{2}}&\lambda&\alpha_1\\ 0 & 1 & \alpha_2\\ 0& 0&\alpha_3 \end{array}\right]
\end{split}
\end{equation}
Using (\ref{weights1}), $V$ can be rewritten as $V=z^T\,z=||z||^2$. On the other hand,
\begin{equation} \label{weights2}
  \begin{split}
\sigma_{min}(\phi \phi^T)||e||^2\leq V=||z||^2 \leq \sigma_{max}(\phi \phi^T)||e||^2
\end{split}
\end{equation}
where $\sigma_{min}(\phi \phi^T)$ and $\sigma_{max}(\phi \phi^T)$ represent the max and min singular values of $\phi \phi^T$ respectively.
Since the Lyapunov function satisfies \eqref{MyEq111} for all $t \in \Re_{\geq0}$, this implies that the whole time during which the adaptation take place is finite. During the finite time, the variables $\hat \theta$ and ${\hat \theta}_d$ cannot escape to infinity since the adaptation laws in \eqref{MyEq99aaa} are well defined. For any bounded conditions $e(0)$, ${\hat \theta}(0)$ and ${\hat \theta}_d(0)$ and by making use of \eqref{MyEq111}, we infer that $e$, ${\hat \theta}$ and ${\hat \theta}_d$ are bounded for all $t \in \Re_{\geq0}$. The proof ends here.
\end{proof}

\end{document}
